\section*{Capítulo \ref{ch:g11:redox} --- Oxidação e redução}

\begin{solution}{exo:g11:redox:1}
\omterm{def:g11:redox:oxidant}{Oxidante}: uma espécie capaz de ganhar \omterm{def:g9:inside-the-atom:nucleus}{elétrons}. \omterm{def:g11:redox:oxidant}{Redutor}: uma espécie
capaz de ceder \omterm{def:g9:inside-the-atom:nucleus}{elétrons}. \omterm{def:g11:redox:oxidation}{Oxidação}: uma perda de \omterm{def:g9:inside-the-atom:nucleus}{elétrons}. \omterm{def:g11:redox:oxidation}{Redução}: um
ganho de \omterm{def:g9:inside-the-atom:nucleus}{elétrons}.
\end{solution}

\begin{solution}{exo:g11:redox:2}
\ce{Fe^{2+}}/\ce{Fe}; \ce{Ag+}/\ce{Ag}; \ce{I2}/\ce{I-};
\ce{Fe^{3+}}/\ce{Fe^{2+}}.
\end{solution}

\begin{solution}{exo:g11:redox:3}
\ce{Ag+ + e- <=> Ag}; \ce{Al^{3+} + 3e- <=> Al};
\ce{Fe^{3+} + e- <=> Fe^{2+}}; \ce{Cl2 + 2e- <=> 2Cl-}.
\end{solution}

\begin{solution}{exo:g11:redox:4}
Uma \omterm{def:g11:redox:oxidation}{oxidação} (\omterm{def:g9:inside-the-atom:nucleus}{elétrons} são perdidos); o ferro age como \omterm{def:g11:redox:oxidant}{redutor}.
\end{solution}

\begin{solution}{exo:g11:redox:5}
O zinco é oxidado e é o \omterm{def:g11:redox:oxidant}{redutor}; o \ce{Cu^{2+}} é reduzido e é o
\omterm{def:g11:redox:oxidant}{oxidante}. Dois \omterm{def:g9:inside-the-atom:nucleus}{elétrons} por \omterm{def:g7:atoms-and-molecules:atom}{átomo} de zinco.
\end{solution}

\begin{solution}{exo:g11:redox:6}
Parte-se de \ce{H2O2 -> H2O}: % ce-unbalanced-ok (step 1 of the method)
oxigênio: dois O à esquerda, logo \ce{2H2O} à direita; hidrogênio: 2 à
esquerda, 4 à direita, logo \ce{2H+} à esquerda; carga: $+2$ à esquerda,
0 à direita, logo dois \omterm{def:g9:inside-the-atom:nucleus}{elétrons} à esquerda:
\ce{H2O2 + 2H+ + 2e- <=> 2H2O}.
\end{solution}

\begin{solution}{exo:g11:redox:7}
\ce{I2 + 2e- -> 2I-} e \ce{2S2O3^{2-} -> S4O6^{2-} + 2e-}; dois
\omterm{def:g9:inside-the-atom:nucleus}{elétrons} em cada uma, logo \ce{I2 + 2S2O3^{2-} -> 2I- + S4O6^{2-}}.
\end{solution}

\begin{solution}{exo:g11:redox:8}
\ce{Cu -> Cu^{2+} + 2e-} e $2\times$ (\ce{Ag+ + e- -> Ag}):
\ce{Cu + 2Ag+ -> Cu^{2+} + 2Ag}. Os \omterm{def:g9:ions:ion}{íons} \ce{Cu^{2+}} formados são azuis.
\end{solution}

\begin{solution}{exo:g11:redox:9}
\ce{C6H8O6 -> C6H6O6 + 2H+ + 2e-} e \ce{I2 + 2e- -> 2I-}:
\ce{C6H8O6 + I2 -> C6H6O6 + 2I- + 2H+}. A vitamina C é o \omterm{def:g11:redox:oxidant}{redutor}.
\end{solution}

\begin{solution}{exo:g11:redox:10}
\ce{Al -> Al^{3+} + 3e-} dá 3 \omterm{def:g9:inside-the-atom:nucleus}{elétrons}, \ce{Cu^{2+} + 2e- -> Cu} toma
2; o menor número comum é 6, por isso a primeira é multiplicada por 2 e
a segunda por 3: \ce{2Al + 3Cu^{2+} -> 2Al^{3+} + 3Cu}.
\end{solution}

\begin{solution}{exo:g11:redox:11}
\ce{H2O2 + 2H+ + 2e- -> 2H2O} e $2\times$ (\ce{Fe^{2+} -> Fe^{3+} + e-}):
\ce{H2O2 + 2H+ + 2Fe^{2+} -> 2H2O + 2Fe^{3+}}.
\end{solution}

\begin{solution}{exo:g11:redox:12}
$n(\ce{Zn}) = 1.30/65.4 = \qty{1.99e-2}{mol}$. Deposita-se um \omterm{def:g7:atoms-and-molecules:atom}{átomo} de
cobre por \omterm{def:g7:atoms-and-molecules:atom}{átomo} de zinco oxidado, logo $n(\ce{Cu}) = \qty{1.99e-2}{mol}$
e $m(\ce{Cu}) = 1.99\times 10^{-2} \times 63.5 = \qty{1.26}{g}$.
\end{solution}

\begin{solution}{exo:g11:redox:13}
\ce{Cr2O7^{2-} + 14H+ + 6e- -> 2Cr^{3+} + 7H2O} e
$6\times$ (\ce{Fe^{2+} -> Fe^{3+} + e-}):
\ce{Cr2O7^{2-} + 14H+ + 6Fe^{2+} -> 2Cr^{3+} + 6Fe^{3+} + 7H2O}.
Seis \ce{Fe^{2+}} por \omterm{def:g9:ions:ion}{íon} dicromato: $6 \times 0.010 = \qty{0.060}{mol}$.
\end{solution}

\begin{solution}{exo:g11:redox:14}
\ce{2ClO- + 4H+ + 2e- -> Cl2 + 2H2O} e \ce{2Cl- -> Cl2 + 2e-}; somando
e dividindo por 2: \ce{ClO- + Cl- + 2H+ -> Cl2 + H2O}. Um ácido fornece
os \omterm{def:g9:ions:ion}{íons} \ce{H+}, e a reação libera gás cloro, um gás tóxico: daí o aviso
para nunca \omterm{def:g2:mixing-and-dissolving:mix}{misturar} os dois produtos.
\end{solution}

\begin{solution}{exo:g11:redox:15}
Os \omterm{def:g9:inside-the-atom:nucleus}{elétrons} não conseguem circular livremente pela solução
(\cref{prop:g11:redox:no-free-electrons}): os que os \omterm{def:g7:atoms-and-molecules:atom}{átomos} de zinco
perdem só podem ser tomados por um \omterm{def:g9:ions:ion}{íon} \ce{Cu^{2+}} que toca o \omterm{def:g9:periodic-table-first-look:metal}{metal}. O
\omterm{def:g7:atoms-and-molecules:atom}{átomo} de cobre se forma, portanto, na superfície da placa, e fica ali.
\end{solution}

\begin{solution}{pb:g11:redox:1}
\textbf{1.} \ce{Cr2O7^{2-}}/\ce{Cr^{3+}} e \ce{CH3COOH}/\ce{CH3CH2OH}.

\textbf{2.} O etanol é oxidado: ele é o \omterm{def:g11:redox:oxidant}{redutor}.

\textbf{3.} \ce{Cr2O7^{2-} + 14H+ + 6e- <=> 2Cr^{3+} + 7H2O}.

\textbf{4.} O carbono está balanceado (dois de cada lado). Oxigênio: dois
à esquerda, um à direita, logo \ce{H2O} à direita. Hidrogênio: quatro à
esquerda, $6 + 2 = 8$ à direita, logo \ce{4H+} à esquerda. Carga: $+4$ à
esquerda, 0 à direita, logo quatro \omterm{def:g9:inside-the-atom:nucleus}{elétrons} à esquerda:
\ce{CH3COOH + 4H+ + 4e- <=> CH3CH2OH + H2O}.

\textbf{5.} Doze \omterm{def:g9:inside-the-atom:nucleus}{elétrons}: $2\times$ a primeira, $3\times$ a segunda
invertida:
\ce{2Cr2O7^{2-} + 3CH3CH2OH + 16H+ -> 4Cr^{3+} + 3CH3COOH + 11H2O}.

\textbf{6.} Os \omterm{def:g9:ions:ion}{íons} \ce{Cr2O7^{2-}} laranja são reduzidos a \omterm{def:g9:ions:ion}{íons}
\ce{Cr^{3+}} verdes onde quer que o etanol chegue até eles.

\textbf{7.} $M = 2\times 39.1 + 2\times 52.0 + 7\times 16.0 =
\qty{294.2}{g/mol}$.

\textbf{8.} $n = 2.0\times 10^{-3}/294.2 = \qty{6.8e-6}{mol}$.

\textbf{9.} Três \omterm{def:g7:atoms-and-molecules:molecule}{moléculas} de etanol para dois \omterm{def:g9:ions:ion}{íons} dicromato:
$n = \tfrac{3}{2}\times 6.8\times 10^{-6} = \qty{1.0e-5}{mol}$.

\textbf{10.} $M(\ce{C2H6O}) = \qty{46.0}{g/mol}$, logo
$m = 1.02\times 10^{-5}\times 46.0 = \qty{4.7e-4}{g}$, cerca de
\qty{0.47}{mg}.

\textbf{11.} $0.25/0.47 \approx 0.53$: cerca de \qty{53}{\percent} do
dicromato.

\textbf{12.} $0.53 \times 10 \approx \qty{5.3}{cm}$.

\textbf{13.} $\qty{0.25}{mg} \times 2100 = \qty{525}{mg}$ por litro de
sangue, cerca de \qty{0.53}{g/L}.

\textbf{14.} O \omterm{def:g4:air-a-mixture-of-gases:air}{ar} expirado encontra primeiro os grãos da entrada; o
dicromato que está ali está em excesso e retira todo o etanol, por isso
a reação é completa num primeiro trecho do tubo, que fica inteiramente
verde, antes que qualquer etanol chegue mais longe.

\textbf{15.} A reação é total: os \omterm{def:g9:ions:ion}{íons} \ce{Cr^{3+}} verdes não voltam a
ser dicromato, por isso um tubo usado não consegue medir de novo.

% ledger: ghs:K2Cr2O7 (H350 among its hazard statements)
\textbf{16.} Ele tem toxicidade aguda (GHS06) e é um perigo grave para a
saúde (GHS08: entre as suas frases de perigo está ``pode provocar
câncer''), além de ser corrosivo e tóxico para os organismos aquáticos;
um aparelho usado pelo público não pode contê-lo.

\textbf{17.} \ce{O2 + 4H+ + 4e- -> 2H2O} e
\ce{CH3CH2OH + H2O -> CH3COOH + 4H+ + 4e-}:
\ce{CH3CH2OH + O2 -> CH3COOH + H2O}.

\textbf{18.} Quatro \omterm{def:g9:inside-the-atom:nucleus}{elétrons}.

\textbf{19.} $n = 0.25\times 10^{-3}/46.0 = \qty{5.4e-6}{mol}$ de
etanol; \omterm{def:g9:inside-the-atom:nucleus}{elétrons} $4n = \qty{2.2e-5}{mol}$, isto é,
$2.2\times 10^{-5}\times 6.02\times 10^{23} = \num{1.3e19}$ \omterm{def:g9:inside-the-atom:nucleus}{elétrons}.

\textbf{20.} $Q = 1.3\times 10^{19}\times 1.60\times 10^{-19} \approx
\qty{2.1}{C}$. Cada \omterm{def:g7:atoms-and-molecules:molecule}{molécula} de etanol manda exatamente quatro \omterm{def:g9:inside-the-atom:nucleus}{elétrons}
pelo fio, por isso a carga é proporcional à quantidade de etanol: medi-la
é medir o etanol.
\end{solution}
